\documentclass[11pt]{scrartcl}
\usepackage{wrapfig}
\usepackage{verbatim} 
\usepackage{acronym}
\usepackage[T1]{fontenc} % Use 8-bit encoding that has 256 glyphs
\usepackage{fourier} % Use the Adobe Utopia font for the document - comment this line to return to the LaTeX default
%\usepackage[english]{babel} % german headers
\usepackage{ucs}             % neccessary for äöü
\usepackage[utf8x]{inputenc} % utf8 support
\usepackage{amsmath,amsfonts,amsthm} % Math packages
\usepackage{hyperref}   % use for hypertext links, including those to external documents and URLs
\usepackage{enumitem}
\usepackage{float}
\usepackage{extramarks} % Required for headers and footers
\usepackage{graphicx} % Required to insert images
\usepackage{subfigure} % Required to insert images
\usepackage{listings} % Required for insertion of code
\usepackage{courier} % Required for the courier font
\usepackage{color}
\usepackage{sectsty} % Allows customizing section commands
\usepackage{epstopdf}
\usepackage[pass]{geometry}
\usepackage{cite}
\usepackage{textcomp}
\usepackage{subfiles}
% Style of citations
\bibliographystyle{plain}  
\usepackage{./sty/fhv}

% Remove indent
\setlength{\parindent}{0pt}

%language
\usepackage[ngerman]{babel}

%general data
\newcommand{\myauthor}{Nicolaj Höss, Marko Petrovi\'c, Kevin Wallis}
\newcommand{\mytitle}{Dokumentation - \textit{fhvRobot}}
\newcommand{\mysubtitle}{für Beaglebone}
\newcommand{\myuniversity}{Fachhochschule Vorarlberg}
\newcommand{\mycourse}{S1: Kopplung und Integration von heterogenen Systemen}
\newcommand{\mystudy}{Master Informatik (ITM3)}
\newcommand{\mydate}{\today{}, Dornbirn}
%document meta data
\title{\mytitle}
\author{\myauthor}
\date{\today{}, Dornbirn}

%header/footer - KOMA
\usepackage{lastpage}
\usepackage[automark]{scrpage2}
\pagestyle{scrheadings}
\clearscrheadfoot
%header left
\ihead[]{\headmark}
%header center
\chead[]{}
%header right
\ohead[]{\includegraphics[width=2cm]{figures/FHV} }
%footer left
\ifoot[]{\myauthor}
%footer center
\cfoot[]{}
%footer right
\ofoot[]{\thepage\ von \pageref{LastPage}}

\usepackage{color}				% defining syntax colours
\definecolor{mygreen}{rgb}{0,0.6,0}
\definecolor{mygray}{rgb}{0.5,0.5,0.5}
\definecolor{mymauve}{rgb}{0.58,0,0.82}

\lstset{ 
  language=Matlab, 
  backgroundcolor=\color{white},
  numbers=left,
  numbersep=5pt,
  showstringspaces=false,
  frame=single,
  rulecolor=\color{black},
  showspaces=false,
  showstringspaces=false,
  showtabs=false,
  tabsize=2,
  breaklines=true,
  captionpos=tl,
  basicstyle=\footnotesize,
  commentstyle=\color{mygreen},
  keywordstyle=\color{blue},     
  stringstyle=\color{mymauve},
  postbreak=\raisebox{0ex}[0ex][0ex]{\ensuremath{\color{red}\hookrightarrow\space}}}


\begin{document}

%%Titelseite
\begin{titlepage}

\begin{minipage}{.5\textwidth} 
  \hspace{12cm}\hfill \includegraphics[width=3cm]{figures/FHV} 
\end{minipage}\vspace{2cm} 

\begin{center}
~
\vfill\vfill\vfill

{\Huge\bfseries\mytitle} \\
{\Large\bfseries\mysubtitle}

\vfill

\Large{eine Arbeit von}

\vfill

{\LARGE\bfseries\myauthor}
\vfill
{\Large\bfseries\mystudy}
\vfill
\vfill
\Large{für die Lehrveranstaltung}

\vfill

{\LARGE\bfseries\mycourse}

\vfill\vfill\vfill
\vfill\vfill\vfill
\vfill\vfill\vfill

\myuniversity
\vfill
\mydate
\end{center}
\end{titlepage}
%%Ende Titelseite

%%Abstract
\begin{abstract}
\thispagestyle{plain}
	
\section*{Kurzfassung}

Diese Arbeit befasst sich mit der Entwicklung eines ...

\end{abstract}

\thispagestyle{plain}
%%Inhaltsverzeichnis
\tableofcontents

\newpage
%%Abbildungsverzeichnis
\listoffigures
\newpage
\pagebreak

\section*{Abkürzungsverzeichnis}

\vspace{1cm}
\begin{acronym}[Bash]
	\acro{OS}{Operating System}
\end{acronym}

\pagebreak 
\section{Allgemein}
In diesem Kapitel werden allgemeine Aspekte zum Robot Projekt erläutert. Dazu zählen insbesondere die durch das Studienprojekt definierten Anforderungen, zusätzlich durch die Studierenden gesetzte Anforderungen und die erreichten Ergebnisse hinsichtlich dieser Anforderungen.

\subsection{Vorgegebene Anforderungen an das Projekt}
Eine Auflistung aller vorgegebenen Anforderungen, insbesondere funktionale Anforderungen an das Betriebssystem sind in Tabelle \ref{table:Prescribed-Requirements} angegeben.

\begin{table}[H]
\begin{tabular}{ p{5cm}| p{9cm} }
  \textbf{Anforderung} & \textbf{Erklärung} \\ 
  \hline
  XX & xxx. \\
  XX & xxx. \\
 \end{tabular}
 \caption{Vorgegebene Anforderungen}
 \label{table:Prescribed-Requirements}
\end{table}

\subsection{Eigene Anforderungen an das Projekt}
Zusätzlich zu den oben angeführten Anforderungen wurden weitere, nicht-funktionale Anforderungen an das Robot Projekt, durch die an der Entwicklung beteiligten Studierenden, definiert. Eine Auflistung aller eigenen Anforderungen sind in Tabelle \ref{table:Own-Requirements} angegeben.

\begin{table}[H]
\begin{tabular}{ p{5cm}| p{9cm} }
  \textbf{Anforderung} & \textbf{Erklärung} \\ 
  \hline
  Hoher Abstraktionsgrad & Alle Komponenten des Projekts sollen einen hohen Abstraktionsgrad aufweisen. \\
  Intuitiver Aufbau & Die Komponenten des Projekts sollen eine intuitive Programmierschnittstelle aufweisen. \\
  Leichte Erweiterbarkeit & Mögliche Erweiterungen sollen ohne große Veränderungen an der Architektur umgesetzt werden können. \\
  Einfache Wartung & Das Robot Projekt soll eine einfache Wartbarkeit hinsichtlich Fehlern aufweisen. \\
 \end{tabular}
 \caption{Vorgegebene Anforderungen}
 \label{table:Own-Requirements}
\end{table}

\subsection{Erfüllung der Anforderungen}
Im Allgemeinen wurden alle zuvor erwähnten funktionalen Anforderungen an das Projekt erfüllt. Einzelne Verbesserungs- bzw. Erweiterungsmöglichkeiten können aus Kapitel \ref{summary} werden. 
\pagebreak 
\section{Projektmanagement}
Im folgenden Abschnitt wird das Projektmanagment und das verwendete Prozessmodell beschrieben. Weiters sind hier die Zugänge zum \textit{Repository} und dem Ticketsystem dokumentiert.

\subsection{Prozessmodell}
Als Prozessmodell wurde \textit{SCRUM}, allerdings mit einigen Adaptionen umgesetzt, wobei das zentrale Vorgehen in Bezug auf Agilität bestmöglich übernommen wurde. Gründe für die Verwendung von \textit{SCRUM} waren vor allem die Möglichkeit zur agilen Umsetzung der vorhandenen und neuer Anforderungen.\\
Es wurde auf das Konzept eines \textit{SCRUM}-Boards verzichtet, stattdessen wurden sämtliche \textit{Stories} als eigenes Ticket in einem passenden Ticketsystem angelegt. Es erfolgte eine Priorisierung der jeweiligen Tickets. Die Abarbeitungsreihenfolge dieser Tickets ergab sich schließlich nach der jeweiligen Priorität selbst.\\

Die angelegten Tickets finden sich unter folgendem Link:

\url{https://github.com/Blackjack92/fhvRobot/issues} \\

Durch Einsicht der offenen und geschlossenen Tickets lässt sich sowohl der Entwicklungsfortschritt, als auch die jeweiligen Designentscheidungen sehr gut nachvollziehen.

\subsection{Versionsverwaltung}
Als Versionsverwaltung wurde \textit{Git} verwendet. Die Entscheidung für die Verwendung von \textit{Git} sind insbesondere die leichte Einbindung im Zusammenhang mit dem angelegten \textit{Repository} und die Möglichkeit der Nicht-linearen Entwicklung.

\subsection{Repository}
\label{Repository}
Das \textit{Repository} für den Source-Code und weitere relevante Dokumente für die Entwicklung des Betriebssystems, wurde auf \textit{Github} angelegt und veröffentlicht. Der Source-Code des Betriebssystems war während der gesamten Entwicklungszeit öffentlich zugänglich. Unter folgendem Link ist das \textit{Repository} einsehbar:\\

\url{https://github.com/Blackjack92/fhvRobot} \\

\subsection{Zeitplan}
Der Zeitplan des Projekts wurde mittels \textit{Microsoft Project} erstellt und verwaltet. Der Zeitplan selbst wurde während des gesamten Projekt im Allgemeinen sehr gut eingehalten. Das angelegte \textit{Microsoft Project}-Projekt ist ebenfalls im oben beschriebenen \textit{Repository} im Ordner \textit{projectmanagement} einsehbar.

\pagebreak 
\section{Architektur}
\label{architecture}
Dieses Kapitel beschreibt die entwickelte Architektur ...

\pagebreak 
\section{Sicherheitsaspekte}
\label{security}

Hinsichtlich der Sicherheit wurden an das ...

\section{Performanzuntersuchungen}
\label{Performanz}

In diesem Kapitel werden Performanzaspekte des ...

\pagebreak 
\section{Zusammenfassung und Ausblick}
\label{summary}

In diesem Kapitel werden die erreichten Ergebnisse zusammengefasst. Weiters wird ein Ausblick auf Möglichkeiten der Weiterentwicklung geboten.

\subsection{Zusammenfassung}

Das primäre Ziel dieses Projektes war ...

\subsection{Ausblick}

...

\subsubsection{Punkte mit Verbesserungspotential}
Einige Punkte des Betriebssystem konnten nicht vollständig abgedeckt werden. Diese Punkte mit Verbesserungspotential werden in Tabelle \ref{table:points-to-improve} kurz beschrieben.

\begin{table}[H]
\begin{tabular}{p{5cm} | p{9cm}}
  \textbf{Sachverhalt} & \textbf{Beschreibung}
  \\ \hline
  XXX & XXX
  \\
 \end{tabular}
 \caption{Übersicht der Punkte mit Verbesserungspotential}
 \label{table:points-to-improve}
\end{table}

\clearpage
\phantomsection
\addcontentsline{toc}{chapter}{Literaturverzeichnis}
\bibliography{./bib/Bibliography}

\end{document}